\subsection{Structural definition and function of the gravitational constant}
\label{sec:ii-definicion-gravedad}

The common realization first receives the length
\(L=(5759/23040)\alpha_{\HMT}^{16}L_*\) from
\eqref{g26:eq:longitud-unica-en}. Its elementary duration is
\(t_0=\pi_{\HMT}L/(54c_{\rm int})\); hence
\(R_{108}=54c_{\rm int}t_0/\pi_{\HMT}=L\).
Equality with the Planck length of this realization,
\(\ell_P=L\), follows after the coupling has been determined.
The positive reference \(L_*\) remains explicit; evaluation in metres
uses the declared section \(L_*=1\,\mathrm m\).

\begin{definition}[Gravitational section contragredient to action]
For a positive section \(\hbar_a\), a velocity
\(c_{\rm int}>0\), and the chronological radius \(R_{108}\) of
section~\ref{sec:gravity}, the circular gravitational section is
\begin{equation}
 G_a=\frac{c_{\rm int}^{3}R_{108}^{2}}{\hbar_a},
 \qquad R_{108}=\frac{54}{\pi}\ell_0,
 \qquad \ell_0=c_{\rm int}t_0.
 \label{eq:ii-definicion-gravedad-seccion}
\end{equation}
It is the action-reciprocal section realizing coincidence
between the cycle's reduced gravitational radius and its
chronological radius, in the convention \(r_g=Gm/c_{\rm int}^2\).
\end{definition}

The radius-coincidence proposition of
section~\ref{sec:gravity} expresses its uniqueness in the collection
\(G_a(\vartheta)=\vartheta c_{\rm int}^{5}t_0^2/\hbar_a\).
The circular condition fixes \(\vartheta=(54/\pi)^2\) and leads
to \eqref{eq:ii-definicion-gravedad-seccion}. That condition and the
positive dimensional sections are part of the realization.
The operator proof in \ref{g26:sec:rigidez-radial-en} extends that
expression to the complete positive character: R4 defines the longitudinal
composition and R5 its radial metric; R8 identifies the resulting operator
with the reduced gravitational radius. With those data, the same coefficient
\(G_{\rm ret}=c_{\rm int}^3L^2/\hbar_{\rm ret}\) acts in every
admissible positive mode. The circular formula is its specialization
to the identity character and the clock conjugate to the length.
Keeping them common to both orientations yields
\begin{equation}
 \frac{G_{\mathrm{pre}}}{G_{\mathrm{ret}}}
 =R_{\mathrm{act}}^{-1},
 \qquad
 \frac{G_a\hbar_a}{c_{\rm int}^{3}}
 =R_{108}^{2}=L^2=\ell_P^2.
 \label{eq:ii-definicion-gravedad-area}
\end{equation}
The proof is direct substitution of
\eqref{eq:ii-definicion-gravedad-seccion}; the action change
is compensated by the reciprocal change in \(G_a\).

This structure assigns gravitation a precise function:
transporting the action section to a common geometric unit.
In \eqref{eq:holst}, \(G\) fixes coupling normalization;
in the area operator, \(G\hbar/c^3\) fixes the area unit.
Reciprocity preserves that unit while the oriented action
reading changes. The reduced state and its occupations subsequently supply
the information and entanglement associated with the boundary.

The tensorial content of the coupling is preserved in
\(\mathbb G_{5,a}=G_aP^{\mathrm{ico}}_5\), constructed in
\ref{sec:ii-pentadica}. Section \(G_a\) determines its common magnitude;
the projector selects the five-component carrier, and \(\Gamma_5\)
transports its orientation to the space of symmetric trace-free tensors.
The basis in \eqref{ii:pent:eq:gravity-helicity2-real}--\eqref{ii:pent:eq:gravity-helicity0-real}
makes the five components explicit. Projection \(\Lambda(n)\) then defines
the transverse reduction. Magnitude, carrier, and reduction are
related operations with different domains within this realization.
